\chapter{The Full Trace: Four Iterations, Every Number Checkable}
\label{ch:full_trace}

Constants, so you can reproduce every line with a pocket calculator:
\begin{align*}
  V(\text{root}) &= 0.97602, \quad c = 1.745, \quad k = 0.33 \\
  P(\text{Kd6}) &= 0.4513, \quad P(\text{Kf6}) = 0.4423, \quad P(\text{Kf5}) = 0.0538, \quad P(\text{Kd5}) = 0.0526 \\[1ex]
  U(a) &= c \cdot P(a) \cdot \frac{\sqrt{N}}{1 + n(a)}, \quad N = \max(1, \text{visits from this node}) \\
  Q_{\text{FPU}} &= Q(\text{parent}) - k \cdot \sqrt{\sum_{\text{visited}} P(a')}
\end{align*}

\section{Iteration 0: Nothing Visited Yet}

Every arrow has $n = 0$, so every $Q$ is the placeholder. Nothing is visited, so the FPU reduction is $k \cdot \sqrt{0} = 0$ and the placeholder is the parent's value itself. $N = \max(1, 0) = 1$.

\begin{center}
\small
\begin{tabular}{lcccc}
  \toprule
  Move $a$ & $n$ & Baseline $Q$ & Curiosity $U = 1.745 \cdot P \cdot 1 / (1+0)$ & Total $S = Q + U$ \\
  \midrule
  \textbf{Kd6} & 0 & 0.97602 & $1.745 \times 0.4513 = 0.78752$ & \textbf{1.76354} \hfill ($\leftarrow \max$) \\
  \textbf{Kf6} & 0 & 0.97602 & $1.745 \times 0.4423 = 0.77181$ & 1.74783 \\
  \textbf{Kf5} & 0 & 0.97602 & $1.745 \times 0.0538 = 0.09388$ & 1.06990 \\
  \textbf{Kd5} & 0 & 0.97602 & $1.745 \times 0.0526 = 0.09179$ & 1.06781 \\
  \bottomrule
\end{tabular}
\end{center}

All four $Q$ are identical, so \textbf{the policy prior alone decides the first move examined.}

\begin{center}
\begin{minipage}{\linewidth}
\centering
% fig_c_iter0.tex -- S-bar strip at Iteration 0
\begin{tikzpicture}
  \node[anchor=west, font=\scriptsize\itshape, text=WorkGrey] at (-5.5, 4.0) {%
    Iteration 0: Initial priors and unvisited FPU baseline ($S = Q_{\mathrm{FPU}} + U$):%
  };
  \vgsbaraxis{-5.5}{5.5}{1.8}
  \vgsbarS{-4.125}{0.97602}{0.78752}{Kd6}{sel}
  \vgsbarS{-1.375}{0.97602}{0.77181}{Kf6}{fpu}
  \vgsbarS{1.375}{0.97602}{0.09388}{Kf5}{fpu}
  \vgsbarS{4.125}{0.97602}{0.09179}{Kd5}{fpu}

  \node[draw=GoldPath, fill=GoldPath!10, rounded corners=2pt, font=\tiny\bfseries, text=GoldPath] at (-4.125, 3.4) {%
    \textbf{SELECTED: Kd6} ($S=1.76354$)%
  };
\end{tikzpicture}

\captionof{figure}{\textbf{Iteration 0 S-bar strip.} Tallest bar is Kd6 ($S = 1.76354$), selected for the first search traversal.}
\label{fig:iter0}
\end{minipage}
\end{center}

\begin{realdata}[Iteration 0 Execution Outcome]
\textbf{Selected: Kd6.} Play it, evaluate the child, $x_1 = +0.96766$ comes back.\\
Now $n(\text{Kd6}) = 1$, $Q(\text{Kd6}) = 0.96766$.\\
The root's own average updates too: $Q(\text{root}) = \text{mean}(0.97602, 0.96766) = 0.97184$.
\end{realdata}

\section{Iteration 1: One Arrow Has a Real Measurement}

Visited policy mass is $0.4513$, so the reduction is $0.33 \cdot \sqrt{0.4513} = 0.33 \cdot 0.67179 = 0.22169$, and unvisited arrows sit at $0.97184 - 0.22169 = 0.75015$.

\begin{center}
\small
\begin{tabular}{lcccc}
  \toprule
  Move $a$ & $n$ & Measured / FPU $Q$ & Curiosity $U = 1.745 \cdot P \cdot 1 / (1+n)$ & Total $S$ \\
  \midrule
  \textbf{Kd6} & 1 & 0.96766 & $1.745 \times 0.4513 / 2 = 0.39376$ & 1.36142 \\
  \textbf{Kf6} & 0 & 0.75015 & $1.745 \times 0.4423 / 1 = 0.77181$ & \textbf{1.52196} \hfill ($\leftarrow \max$) \\
  \textbf{Kf5} & 0 & 0.75015 & $1.745 \times 0.0538 / 1 = 0.09388$ & 0.84403 \\
  \textbf{Kd5} & 0 & 0.75015 & $1.745 \times 0.0526 / 1 = 0.09179$ & 0.84194 \\
  \bottomrule
\end{tabular}
\end{center}

\texttt{Kd6}'s evaluation came back \textbf{excellent} ($+0.96766$) --- and its score still \textit{fell}, from $1.76354$ to $1.36142$. Two separate things did that:
\begin{enumerate}[leftmargin=2em]
  \item Its $U$ halved, because $n$ went $0 \to 1$.
  \item Its $Q$ dropped slightly, the borrowed $0.97602$ giving way to the measured $0.96766$.
\end{enumerate}
Meanwhile \texttt{Kf6} kept its full curiosity bonus.

\begin{center}
\begin{minipage}{\linewidth}
\centering
% fig_c_iter1.tex -- S-bar strip at Iteration 1
\begin{tikzpicture}
  \node[anchor=west, font=\scriptsize\itshape, text=WorkGrey] at (-5.5, 4.0) {%
    Iteration 1: Kd6 measured ($Q{=}0.96766$, $U$ halved); unvisited FPU drops to $0.75015$:%
  };
  \vgsbaraxis{-5.5}{5.5}{1.8}
  \vgsbarS{-4.125}{0.96766}{0.39376}{Kd6}{meas}
  \vgsbarS{-1.375}{0.75015}{0.77181}{Kf6}{sel}
  \vgsbarS{1.375}{0.75015}{0.09388}{Kf5}{fpu}
  \vgsbarS{4.125}{0.75015}{0.09179}{Kd5}{fpu}

  \node[draw=GoldPath, fill=GoldPath!10, rounded corners=2pt, font=\tiny\bfseries, text=GoldPath] at (-1.375, 3.4) {%
    \textbf{SELECTED: Kf6} ($S=1.52196$)%
  };
\end{tikzpicture}

\captionof{figure}{\textbf{Iteration 1 S-bar strip.} Kf6 takes the lead ($S = 1.52196$) because Kd6's curiosity bonus halved.}
\label{fig:iter1}
\end{minipage}
\end{center}

\begin{realdata}[Iteration 1 Execution Outcome]
\textbf{Selected: Kf6.} $x_1 = +0.98598$ comes back --- even better than Kd6's.\\
$Q(\text{root}) = \text{mean}(0.97602, 0.96766, 0.98598) = 0.97655$.
\end{realdata}

\section{Iteration 2: Both Good Moves Now Measured}

Visited policy mass is $0.4513 + 0.4423 = 0.8936$, so the reduction grows to $0.33 \cdot \sqrt{0.8936} = 0.33 \cdot 0.94531 = 0.31195$, and unvisited arrows drop to $0.97655 - 0.31195 = 0.66460$.

$N = 2$ now, so $\sqrt{N} = 1.41421$.

\begin{center}
\small
\begin{tabular}{lcccc}
  \toprule
  Move $a$ & $n$ & Measured / FPU $Q$ & Curiosity $U = 1.745 \cdot P \cdot 1.41421 / (1+n)$ & Total $S$ \\
  \midrule
  \textbf{Kd6} & 1 & 0.96766 & $1.745 \times 0.4513 \times 1.41421 / 2 = 0.55696$ & 1.52462 \\
  \textbf{Kf6} & 1 & 0.98598 & $1.745 \times 0.4423 \times 1.41421 / 2 = 0.54585$ & \textbf{1.53183} \hfill ($\leftarrow \max$) \\
  \textbf{Kf5} & 0 & 0.66460 & $1.745 \times 0.0538 \times 1.41421 / 1 = 0.13279$ & 0.79739 \\
  \textbf{Kd5} & 0 & 0.66460 & $1.745 \times 0.0526 \times 1.41421 / 1 = 0.12983$ & 0.79443 \\
  \bottomrule
\end{tabular}
\end{center}

\begin{center}
\begin{minipage}{\linewidth}
\centering
% fig_c_iter2.tex -- S-bar strip at Iteration 2
\begin{tikzpicture}
  \node[anchor=west, font=\scriptsize\itshape, text=WorkGrey] at (-5.5, 4.0) {%
    Iteration 2: Both winning moves measured ($N=2, \sqrt{N}=1.414$); Kf6 narrowly edges Kd6:%
  };
  \vgsbaraxis{-5.5}{5.5}{1.8}
  \vgsbarS{-4.125}{0.96766}{0.55696}{Kd6}{meas}
  \vgsbarS{-1.375}{0.98598}{0.54585}{Kf6}{sel}
  \vgsbarS{1.375}{0.66460}{0.13279}{Kf5}{fpu}
  \vgsbarS{4.125}{0.66460}{0.12983}{Kd5}{fpu}

  \node[draw=GoldPath, fill=GoldPath!10, rounded corners=2pt, font=\tiny\bfseries, text=GoldPath] at (-1.375, 3.4) {%
    \textbf{SELECTED: Kf6 (2nd time)} ($S=1.53183$)%
  };
\end{tikzpicture}

\captionof{figure}{\textbf{Iteration 2 S-bar strip.} Both winning moves measured; Kf6 ($S = 1.53183$) edges Kd6 ($S = 1.52462$) on quality.}
\label{fig:iter2}
\end{minipage}
\end{center}

\begin{realdata}[Iteration 2 Execution Outcome]
\textbf{Selected: Kf6 --- for the second time.}
\end{realdata}

Note how close it was: $1.53183$ against $1.52462$. \texttt{Kf6} won only because its measured value came back higher. The two bad moves have fallen to about $0.79$ and are no longer competitive.

\section{Iteration 3: And This Is the Answer to ``Why Revisit a Move?''}

\texttt{Kf6} is entered again. Its second sample comes back lower, $x_2 = +0.95128$, and the average moves accordingly:
\begin{equation}
  Q(\text{Kf6}) = \frac{0.98598 + 0.95128}{2} = \mathbf{0.96863} \quad (n(\text{Kf6}) = 2)
\end{equation}

Check with the incremental form:
\begin{equation*}
  0.98598 + \frac{0.95128 - 0.98598}{2} = 0.98598 - 0.01735 = \mathbf{0.96863} \quad \text{(Agrees)}
\end{equation*}

The unvisited arrows drop again, to $0.97024 - 0.31195 = 0.65829$, because the root's own running average moved.

\begin{center}
\begin{minipage}{\linewidth}
\centering
% fig_c_iter3.tex -- S-bar strip before Iteration 3 pick
\begin{tikzpicture}
  \node[anchor=west, font=\scriptsize\itshape, text=WorkGrey] at (-5.5, 4.0) {%
    Scores compared BEFORE Iteration 3 pick ($S = Q + U$):%
  };
  \vgsbaraxis{-5.5}{5.5}{1.8}
  \vgsbarS{-4.125}{0.96766}{0.55696}{Kd6}{meas}
  \vgsbarS{-1.375}{0.98598}{0.54585}{Kf6}{sel}
  \vgsbarS{1.375}{0.66460}{0.13279}{Kf5}{fpu}
  \vgsbarS{4.125}{0.66460}{0.12983}{Kd5}{fpu}

  \node[draw=GoldPath, fill=GoldPath!10, rounded corners=2pt, font=\tiny\bfseries, text=GoldPath] at (-1.375, 3.4) {%
    \textbf{SELECTED: Kf6 (2nd time)} ($S=1.53183$)%
  };
\end{tikzpicture}

\captionof{figure}{\textbf{Iteration 3 S-bar strip.} Pre-pick comparison that led to Kf6's second visit.}
\label{fig:iter3}
\end{minipage}
\end{center}

\textbf{The important part is not the arithmetic --- it is where that second sample came from.} See Chapter~\ref{ch:three_questions}.
